\documentclass[conference]{IEEEtran}

\usepackage{graphicx} 
\usepackage{amsmath}
\usepackage{amssymb} 
\usepackage{amsthm}
\usepackage{cite}
\usepackage{url}
\usepackage{multirow}

\title{NeuroCursor: A Behavioral Biometric Framework for Mouse-Based User Authentication}

\author{
\IEEEauthorblockN{
Shaurya Saxena,
Chetan Gangireddy,
Shreehan Aalok Pathak,
Hamid Abdul Mossa
}
\IEEEauthorblockA{
%Date
Computer Science Department\\
The University of Texas at Dallas, Richardson, Texas\\
Email: Shaurya.saxena090@gmail.com
%ORCID
}
}


\begin{document}

\maketitle

\begin{abstract}
    Behavioral biometrics offer a promising alternative to conventional authentication methods and processes by identifying users through patterns with respect to their interactions with digital devices and software. This paper presents NeuroCursor, a mouse-dynamics-based biometric authentication system that analyzes the subtle variations in cursor movement over a certain duration where a user engages in an iterated point-and-click task. During enrollment, users complete a sequence of interactions with cursor targets appearing at arbitrary screen locations. The system records behavioral features such as cursor trajectory, movement duration, velocity, acceleration, path efficiency, curvature, click timing, and directional changes. These features are processed and analyzed using machine learning models to differentiate between legitimate users and potential imposters. 

    The proposed system is designed to evaluate whether a small number of cursor interactions can provide an adequate amount of behavioral information for reliable identity verification while simultaneously maintaining a simple and minimally intrusive user experience. We examine the effectiveness of different sets of features and classification models using various metrics, such as accuracy, precision, recall, F1-score, false acceptance rate, and false rejection rate. Preliminary results indicate that mouse movement patterns contain measurable user-specific characteristics, likely suggesting that NeuroCursor may serve as an additional authentication layer for desktop applications and other cursor-based systems with a purpose similar to multi-factor authentication. The study also discusses challenges involving behavioral variability, data collection processes, model generalization, and resistance to imitation and overfitting. Ultimately, NeuroCursor demonstrates the potential implications of machine-learning-based mouse dynamics as an accessible form of behavioral biometric authentication. 
\end{abstract}

\begin{IEEEkeywords}
behavioral biometrics, mouse dynamics, cursor trajectory analysis, user authentication, machine learning, human--computer interaction, cybersecurity, biometric security, feature extraction, Convolution Neural Network (CNN), Gated Recurrent Unit (GRU), spatiotemporal feature learning, deep learning 
\end{IEEEkeywords}

\section{Introduction}
\subsection{Background and Motivation}
Authentication systems are fundamental to protecting digital accounts, devices, and information systems. Conventional authentication customarily relies on knowledge-based credentials, such as passwords and personal identification codes, or possession-based credentials for validation, such as security tokens and mobile devices. Albeit these mechanisms may be profound, they primarily verify that a user possesses the correct credential rather than verifying whether the person interacting with a given system is the legitimate account holder. Credentials may also be forgotten, shared, investigated, or compromised, facilitating the development of authentication methods that incorporate measurable characteristics unique to the user. 

Biometric authentication recognizes or validates individuals utilizing physiological or behavioral characteristics. Physiological biometrics--including fingerprints, facial features, and iris patterns--are associated with the physical attributes of an individual. Behavioral characteristics, by contrast, characterize patterns in how a specific individual performs an act or interacts with a system in an environment. Examples of such behavioral characteristics include keystroke timing and duration, touchscreen gestures, gait, signature dynamics, and mouse movement behavior\cite{nistBiometrics, khanMouseSurvey}. Because behavioral signals can oftentimes be obtained through existing input devices, they may provide an additional authentication method without the need for further specialized biometric hardware. 

Mouse dynamics is a behavioral biometric modality that is based on measurable properties and metrics of cursor interaction between a user and a software environment. A cursor trajectory contains more information than its mere starting and ending positions. It may also encode movement duration, changes in direction, path length, velocity, acceleration, curvature, pauses, corrective motion, overshoot, and click timing. Such characteristics are influenced by an individual's motor control, reaction patterns, hand-eye coordination, and interaction habits. Even though no single mouse is expected to unique identify a person, a series of movements may contain patterns that are statistically significant for distinguishing genuine users from imports\cite{khanMouseSurvey,deutschmannContinuousAuthentication}.

This paper presents \textit{NeuroCursor}, a behavioral biometric authentication and validation system based on a randomized point-and-click task. During an authentication attempt, targets are displayed at varying positions within the user interface of the platform. Then, the user moves the cursor to each target and selects it, while the system records the resulting trajectory as a multivariate time series. Randomization changes the required movement direction and distance between the consecutive interactions, thereby encouraging the model to learn characteristics of how a user moves rather than merely memorizing one fixed trajectory and overfitting. 

Each recorded trajectory may be represented as 

\begin{equation}
\mathbf{X}
=
\left[
\mathbf{x}_1,
\mathbf{x}_2,
\ldots,
\mathbf{x}_T
\right]^{\top}
\in
\mathbb{R}^{T \times d},
\label{eq:trajectory matrix}
\end{equation}

where $T$ denotes the number of sampled time steps, $d$ denotes the number of recorded or derived variables, and 

\begin{equation}
\mathbf{x}_T = 
\begin{bmatrix}
x_t &
y_t &
\Delta t_t & 
v_{x,t} & 
v_{y,t} & 
a_{x,t} & 
a_{y,t}
\end{bmatrix}^{\top}
\label{eq:trajectory_vector}
\end{equation}

is one possible feature representation at time $t$. Additional variables--such as jerk, curvature, movement angle, button state, and distance from target--may be incorporated depending on the final implementation. 

NeuroCursor uses a hybrid Convolutional Neural Network (CNN) and Gated Recurrent Unit (GRU) architecture. The convolution component is intended to learn short-range movement structures, including local velocity changes, micro-corrections, and deceleration patterns. The GRU component processes the order convolution representations to model dependencies across the recorded trajectories. This division is motivated by prior research pertaining to time series, which shows that convolutional and recurrent components can provide complementary forms of representation \cite{foumaniTimeSeriesSurvey, elsayedGRUFCN}.

Fundamentally, the intended purpose of NeuroCursor is not necessarily to replace passwords or established authentication methods and processes. Rather, mouse behavior may function as an additional validation signal in a layered security framework. For instance, a low-confidence behavioral score could trigger an additional authentication process, whereas a high confidence signal could support a less obtrusive login experience for users. Such an approach recognizes that behavioral biometrics are probabilistic and may vary with fatigue, injury, hardware, the local environment, or changes in user behavior.   

\subsection{Problem Statement}

The central problem addressed in this work is whether a short sequence of randomized cursor movements contains the sufficient user-specific information to support and validate biometric verification. In contrast to identification, validation tests whether an observed sample is consistent with a claimed category--the respective account for the intended user. 

Let $u$ denote a claimed user and let $\mathbf{X}$ denote an observed cursor sequence. the authentication model estimates 

\begin{equation}
s(\mathbf{X},u)
=
P(Y=1 \mid \mathbf{X}, u), 
\label{eq:auth_probability}
\end{equation}

where $Y=1$ indicates that the sequence was generated by the claimed anticipated user, and $Y=0$ indicates that it was generated by an imposter. Given an authentication threshold $\tau$, the predicted decision is 

\begin{equation}
\widehat{Y}
=
\begin{cases}
1, & s(\mathbf{X}, u) \geq \tau,\\
0, & s(\mathbf{X}, u) < \tau.
\end{cases}
\label{eq:threshold_decision}
\end{equation}

However, the threshold introduces a tradeoff in terms of security and utility. Increasing $\tau$ generally results the system being more conservative, which potentially reduces unauthorized acceptances while simultaneously increasing the number of legitimate users who are rejected. Decreasing $\tau$ may allow for augmented convenience, but it also increases security risk and vulnerability. Consequently, NeuroCursor may be evaluated using biometric error measure such as the false acceptance rate and false rejection rate, rather than merely relying exclusively on classification accuracy. 

For an authentication attempt containing $K$ point-and-click movements, let

\begin{equation}
\mathcal{Q}
=
\left\{
\mathbf{X}^{(1)},
\mathbf{X}^{(2)},
\ldots,
\mathbf{X}^{(K)}
\right\}
\label{eq:challenge_set}
\end{equation}

represent the complete challenge sequence. An aggregate authentication score may be defined as 

\begin{equation}
S(\mathcal{Q},u)
=
\frac{1}{K}
\sum_{k=1}^K
s\left(\mathbf{X}^{(k)},u\right).
\label{eq:aggregate_score}
\end{equation}

For the proposed NeuroCursor interaction, the intended value is $K=5$. Nevertheless, distinct experiments comparing different values of $K$ can determine whether incorporating additional movements yields a statistically significant improvement in the authentication method's reliability or merely prolongs the interaction time, thereby curtailing the overall efficiency of the authentication process. 

Several technical challenges arise. Cursor trajectories vary in terms of duration and number of recorded samples, hence requiring resampling, padding, masking, or another methodology for handling unequal sequence lengths. Movement behavior can also be impacted by mouse sensitivity, DPI, acceleration settings, screen resolution, target geometry, input hardware, and the collection environment. Recent empirical research has corroborated that mouse configuration parameters may substantially alter prevalent mouse-dynamics features, resulting in device control and normalization being crucial for credible evaluation \cite{kuricMouseCredibility}. The restrained size of a student-collected dataset may additionally amplify the risk of potentially overfitting, especially for a deep neural network. 

The research problem can therefore be summarized as the design of a model that learns user-discriminative patterns while constraining the model's sensitivity to irrelevant variation in such patterns. The system must discern between variation that reflects identity and variation caused by target placement, sequence length, sampling rate, device settings, and transient behavioral changes.

\subsection{Research Questions and Hypotheses}
This study is guided by the following research questions:

\begin{enumerate}
    \item Does the proposed CNN--GRU model outperform traditional machine-learning models and single-component neural architectures that exist?

    \item How does the authentication model's performance change when one, three, or five cursor movements are agglomerated into a single decision collectively?

    \item Which spatial, temporal, and kinematic features contribute the most strongly to user authentication and validation processes?

    \item How does the authentication threshold impact the false acceptance rate, false rejection rate, and equal error rate?

    \item To what extent does the model's performance remain stable across discrete collection trials and user sessions?

    \item How sensitive is the model to variation in target distance, movement direction, hardware conditions, and sampling characteristics?
\end{enumerate}

The primary statistical hypothesis concerns the separation between genuine and imposter authentication scores. Let $S_G$ and $S_I$ denote the random variable representing scores from genuine attempts and the corresponding variable for imposter attempts, respectively. The null and alternative hypotheses are 

\begin{equation}
H_0:
\mathbb{E}[S_G],
\leq
\mathbb{E}[S_I]
\label{eq:null_hypothesis}
\end{equation}

\begin{equation}
H_1:
\mathbb{E}[S_G],
>
\mathbb{E}[S_I]
\label{eq:alternative_hypothesis}
\end{equation}

Failure to accept $H_0$ would suggest that the model yields higher scores to genuine samples systematically. Despite this, the statistical separation alone would not evince the pragmatic effectiveness of such an authentication method. The magnitude of the separation, confidence intervals, effect size, and resulting biometric error rates must also be taken into consideration when evaluating the biometric system. 

The following model-level hypotheses are additionally proposed:

\begin{itemize}
\item \textbf{H1:} The CNN--GRU architecture will achieve a lower equal rate than CNN-only and GRU-only models due to the combination of local motion-feature extraction with sequential modeling. 

\item \textbf{H2:} Aggregating five target interactions will yield a lower false acceptance rate and equal error rate in comparison to the use of one interaction. 

\item \textbf{H3:} The usage of derived kinematic features--such as velocity, acceleration, curvature, jerk, and path efficiency--will demonstrate better performance relative to raw cursor coordinates alone. 

\item \textbf{H4:} Session-separated model validation will result in lower performance than a random sample-level split because it more accurately measures generalization to new authentication attempts.
\end{itemize}

However, these hypotheses should only be retained if proven by the experiments/ Particularly, any statements regarding the superiority of the proposed hybrid model should be confirmed with proper baselines and ablation experiments. 

\subsection{Contributions}
\begin{itemize}
    \item We introduce NeuroCursor, a randomized point-and-click system for garnering mouse dynamics authentication sequences. 

    \item We represent cursor behavior as a multivariate time series that encompasses both raw positional data and higher-order features--such as spatial, temporal, geometric, and kinematics. 

    \item We propose a hybrid CNN--GRU model where the convolution layers are used to acquire localized movement patterns and the recurrent layers learn the temporal structure of such patterns. 

    \item We evaluate the system as a biometric verification problem using a false acceptance rate, false rejection rate, receiver operating characteristic analysis, area under the curve, and equal error rate in addition to conventional classification metrics. 

    \item We examine the relationship between the number of cursor interactions, interaction accuracy, and interaction latency. 

    \item We compare the complete model to classical machine-learning baselines--such as CNN-only and GRU-only architectures, as well as limited feature configurations--to determine the value of each component to informative signals. 
\end{itemize}

\section{Related Work}
\subsection{Behavioral Biometrics}

Biometric technologies include physiological and behavioral types. Physiological biometrics are used to acquire identity information from physical attributes, whereas behavioral biometrics depend on patterns in actions or interactions. According to NIST, a biometric refers to any measurable physiological or behavioral characteristics that can be used to identify or validate a claimed identity \cite{nistBiometrics}. It should be noted that this definition implies that mouse dynamics is placed within the broader field of biometric authentication while underscoring the distinction between identification and verification. 

Various behavioral biometric systems have been developed using keystroke dynamics, handwriting, touchscreen gestures, gait, voice behavior, and patterns of computer interaction. The principal advantage of using such behavioral signals is that numerous behavioral signals can be captured through devices that are already used during primary interactions, allowing passive or minimally invasive verification. However, behavioral characteristics are not perfectly invariant; user behavior varies depending on physical condition, emotional state, familiarity with tasks, posture, device type, and environmental conditions. Therefore, behavioral biometric matching is generally probabilistic in nature rather than deterministic. 

The behavioral authentication process can be performed either as a discrete checkpoint or continuously after the initial login. Discrete authentication evaluates a specific interaction created for verification, whereas continuous authentication monitors behavior during an active session. Deutschmann et al. studied continuous behavioral authentication using keystroke and mouse data collected from ninety-nine participants over a ten-week period, exemplifying both the potential and longitudinal difficulty of behaviometric verification \cite{deutschmannContinuousAuthentication}. Continuous authentication methods may help detect account takeover after login, but they oftentimes require a prolonged duration of time and involve additional privacy risks related to the ongoing collection of user activity throughout sessions. 

NeuroCursor employs challenge-based authentication rather than unrestrained, continuous monitoring. Users intentionally execute a short sequence of randomized point-and-click actions at the authentication stage. This overall structure allows for  greater control over task conditions and a clear beginning and end for each sample. Moreover, this methodology also communicates the authentication intent more directly rather than passively monitoring ordinary cursor activity and trajectory. 

\subsection{Mouse-Dynamics Authentication}

Mouse dynamics refers to the analysis of cursor trajectories and mouse-button interactions for the ultimate purpose of characterizing user behavior. Recorded characteristics may include, for instance, cursor coordinates, timestamps, movement direction, distance, velocity, acceleration, curvature, pause duration, clicks, double-clicks, drag events, and interactions with graphical interface elements. Features can be extracted either from an individual user's actions or from a longer sequences of activity. 

Task design in existing mouse-dynamics research varies substantially. Some systems, for instance, acquire unconstrained cursor activity while a user performs typical computer tasks. Other systems use constrained tasks--such as target selection, drag operation, object trace drawing, or interaction with predefined interface widgets. Although unconstrained collection may be more relevant to the representation of normal behavior in users, differences in the content of applications and tasks can augment variability. Constrained tasks improve comparability but may be more intrusive and less relevant to natural user experiences. 

Khan \textit{et al.} surveyed mouse dynamics and closely related widget interaction, classifying prior work by data collection tasks, feature definitions, datasets, algorithms, fusion methods, and limitations \cite{khanMouseSurvey}. Their review demonstrates that the field has utilized various statistical techniques, classical machine learning classifiers, and deep-learning algorithms. It also identifies tenacious challenges involving discrepant experimental protocols, a paucity of publicly available data, hardware variance, and difficulty comparing results across studies. 

Geometrically, the cursor trajectory is a sampled planar curve. Given cursor position 

\begin{equation}
\mathbf{r}(t)
=
\begin{bmatrix}
x(t)\\
y(t)
\end{bmatrix},
\end{equation}

its velocity and acceleration vectors are 

\begin{equation}
\mathbf{v}(t)
=
\frac{d\mathbf{r}(t)}{dt},
\qquad
\mathbf{a}(t)
=
\frac{d^2\mathbf{r}(t)}{dt^2}.
\end{equation}

In terms of the continuous formulation, the cursor's trajectory length is defined as

\begin{equation}
L
=
\int_{t_0}^{t_f}
\left\|
\frac{d\mathbf{r}(t)}{dt}
\right\|_2
\,dt.
\label{eq:continuous_path_length}
\end{equation}

Since NeuroCursor registers cursor positions at distinct, discrete timestamps, the trajectory length is computed numerically as 

\begin{equation}
L=
\sum_{i=2}^{T}
\left\|
\mathbf{r}_i-\mathbf{r}_{i-1}
\right\|_2.
\label{discrete_path_length}
\end{equation}

The path efficiency between starting position $\mathbf{r}(t_0)$ and the target $\mathbf{q}$ can be expressed as 

\begin{equation}
\eta
=
\frac{
\left\|
\mathbf{q}-\mathbf{r}(t_0)
\right\|_2
}{
L
},
\qquad
0 < \eta \leq 1.
\label{eq:path_efficiency}
\end{equation}

These quantities capture different aspects and properties of user behavior. Two different users may reach the same target point in similar amounts of time but differ in path curvature, acceleration profile, corrective motion, or deceleration near the target element. 

One major concern about the methodology pertaining to measuring the mouse-dynamics is the question of whether or not the measured behaviors are due to either the user or the hardware itself. Kuric \textit{et al.} investigated the dependence of mouse-dynamics variables on mouse configuration parameter configurations and discovered that commonly analyzed mouse-dynamics features can be substantially impacted by configuration differences \cite{kuricMouseCredibility}. Accordingly, an authentication performance on one fixed system may not generalize to another mouse, trackpad, sensitivity setting, or screen configuration. This concern motivates the recording of experimental circumstances and conditions, normalizing spatial data, and distinguishing claims about same-device verification from claims regarding cross-device verification. 

NeuroCursor differs from long-window monitoring systems in the sense that it acquires and aggregates a small, predefined number of directed movement. For further elaboration, NeuroCursor differs from a fixed gesture as target points in NeuroCursor are randomly generated. Consequently, the model receives multiple movement directions and distances while the interaction remains brief enough to function as an authentication challenge. 

\subsection{Deep Learning for Time-Series Biometrics}

\section*{Acknowledgment}
The authors appreciate Dr. Anurag Nagar for his guidance and mentorship throughout this research process. The authors also thank Hamid Rouhani, a concurrent Ph.D. student in computer science at The University of Texas at Dallas, for his valuable technical guidance, research support, feedback, and contributions to the development and evaluation of NeuroCursor. 

\end{document}
